\section{Relative bounds and stability}\label{sec:stability}

The preceding sections provide two estimates with different roles.
The zero-mode bound controls a local observable after particles have
been removed. The loss evolution connects an arbitrary state to that
zero space at a cost measured by its initial interaction energy.
We first combine them into a relative form bound. Spectral transport
will then put the disorder perturbation in the class covered by that
bound.

Throughout this section, $q=3(N-1)$ is sufficiently large,
$\Omega\in Z_{N,q}$ is the unit Laughlin vector, and $H=H_q$.
Choose the fixed $\mu>0$ from
Proposition~\ref{prop:grand-canonical-coercivity} and write
\begin{equation}\label{eq:stability-G}
 G=H+\mu(N-\mathcal N).
\end{equation}
Thus $G$ is positive, has kernel $\mathbb C\Omega$, and has a uniform
gap. Moreover it controls both $H$ and $|N-\mathcal N|$.

\subsection{A relative bound for terms that annihilate the ground state}

We use the local norm $\|A\|_{x,\ell}$ of
\eqref{locality:ball-norm} at one fixed order $\ell$.
Partition the sphere into cells of uniformly bounded diameter whose
centers satisfy the two-dimensional volume bound of
Section~\ref{sec:locality}. For a continuous interaction, collect
all terms centered in a cell and assign their integral to that cell's
center. Moving a center a bounded distance changes the local norm by
a fixed factor. We normalize the total localization norm in each cell
to at most one.

Physical operators act trivially on the complementary slot modes;
the inequalities below are on physical Fock space. Averaging ball
approximants over total-slot-number rotations preserves support and
approximation error, so approximants of a neutral operator may be
chosen neutral. Taking Hermitian parts gives Hermitian approximants.

Choose $\ell$ larger than the polynomial support-radius exponent in
Theorem~\ref{thm:charge-local-linear}, and larger than any fixed
geometric exponents needed below. A dyadic decomposition into ball
approximants shows that the zero-mode estimate extends to interactions
with bounded total $\|\cdot\|_{x,\ell}$ per cell. Explicitly, the
shell of radius $2^k$ has norm $O(2^{-k\ell})$, while its zero-mode
cost is $O(2^{kB})$ for a fixed $B<\ell$; the series converges.

\begin{proposition}[Relative form bound]\label{prop:relative-bound}
For this fixed sufficiently large $\ell$, there is a constant $C$
independent of $q$ such that
\begin{equation}\label{eq:relative-bound}
 -CG\ \le\ \sum_x A_x\ \le\ CG
\end{equation}
whenever the $A_x$ are physical, neutral, Hermitian operators satisfying
$A_x\Omega=0$ and $\|A_x\|_{x,\ell}\le1$.
If the sum or integral of localization norms in each cell is at most
$S$, the conclusion becomes $-CSG\le\sum_xA_x\le CSG$.
\end{proposition}

\begin{proof}
Fix $q$ for the moment and let $K=K_q$ be the least constant in
\eqref{eq:relative-bound} for the stated class. It is finite: every
admissible sum annihilates $\Omega$ on both sides, its norm is bounded
by the finite number of cells, and $G$ has a positive gap on
$\Omega^\perp$. This observation gives no useful uniform bound;
we obtain one from the loss evolution.

Let $J_y$ be the jumps of Theorem~\ref{thm:local-energy-descent},
including their bounded auxiliary indices in the measure $dy$.
They obey
\[
 \int J_y^*J_y\,dy=H,
 \qquad J_y\Omega=0\quad\text{for almost every }y.
\]
The second assertion follows by taking the expectation of the first
in $\Omega$. Put $C_{xy}=[A_x,J_y]$ and
$D_{xy}=1+d(x,y)$. Since $A_x$ is even, the ordinary commutator
vanishes for disjointly supported approximants, including when $J_y$
is odd. Lemma~\ref{locality:calculus}(iii) gives
\begin{equation}\label{eq:commutator-norms}
 \|C_{xy}\|\le C D_{xy}^{-\ell},
 \qquad \|C_{xy}\|_{x,\ell}\le C,
 \qquad
 \|C_{xy}^*C_{xy}\|_{x,\ell}\le C D_{xy}^{-\ell}.
\end{equation}
The last estimate retains the same order $\ell$ required of $A_x$.
Thus, after integration in $y$, the positive operators
$C_{xy}^*C_{xy}$ remain in the class whose best relative-bound constant
is $K$. Fixing one polynomial norm makes this closure possible.

Every $C_{xy}$ annihilates $\Omega$, because both $A_x$ and $J_y$
do. Therefore $C_{xy}^*C_{xy}$ is a physical, positive, neutral,
centered operator to which the definition of $K$ applies. Choose
\begin{equation}\label{eq:distance-weight-choice}
 2<b<\ell-2.
\end{equation}
The two-dimensional volume bound and \eqref{eq:commutator-norms}
give
\begin{equation}\label{eq:squared-commutator-sum}
 \sum_x\int D_{xy}^{b}C_{xy}^*C_{xy}\,dy\le C K G.
\end{equation}
Indeed for each fixed $x$ the integrated operator has localization
norm at most $C\int D_{xy}^{b-\ell}\,dy\le C$, and kills
$\Omega$. Collect it as a single admissible term at cell $x$.
This also explains the continuous-family extension in the statement.

Let $\rho_t$ be the state evolution dual to the loss generator,
starting at an arbitrary density matrix $\rho_0$. All observables in
the desired inequality preserve number, so replacing $\rho_0$ by
its number-diagonal part changes none of the relevant expectations.
For $A=\sum_xA_x$, the Lindblad identity reads
\[
 \frac{d}{dt}\operatorname{Tr}(\rho_t A)
 =\operatorname{Re}\sum_x\int
       \operatorname{Tr}(\rho_t J_y^*C_{xy})\,dy.
\]
Apply Cauchy--Schwarz to this sum with weights $D_{xy}^{-b}$ and
$D_{xy}^{b}$. Since $\sum_xD_{xy}^{-b}\le C$ uniformly in $y$,
\eqref{eq:squared-commutator-sum} yields
\begin{equation}\label{eq:relative-derivative}
 \left|\frac{d}{dt}\operatorname{Tr}(\rho_t A)\right|
 \le C\bigl(\operatorname{Tr}(\rho_t H)\bigr)^{1/2}
          \bigl(K\operatorname{Tr}(\rho_t G)\bigr)^{1/2}.
\end{equation}
The energy-descent estimates give constants $c,C>0$ such that
\[
 \operatorname{Tr}(\rho_t H)\le e^{-ct}
          \operatorname{Tr}(\rho_0 H),\qquad
 \operatorname{Tr}(\rho_t G)\le C\operatorname{Tr}(\rho_0 G),
 \qquad H\le CG.
\]
Integrating \eqref{eq:relative-derivative} therefore costs at most
$C\sqrt K\operatorname{Tr}(\rho_0 G)$.

Along a sequence $t_j\to\infty$, finite dimensionality gives a
limit state $\rho_\infty$ supported on $Z_q$. Its sectors have
$n\le N$ by Lemma~\ref{lem:zero-modes}. The zero-mode bound and
$\langle\Omega,A\Omega\rangle=0$ show that
\[
 |\operatorname{Tr}(\rho_\infty A)|
 \le C\operatorname{Tr}(\rho_\infty(N-\mathcal N))
 =C\mu^{-1}\operatorname{Tr}(\rho_\infty G)
 \le C\operatorname{Tr}(\rho_0G).
\]
Combining the two bounds, and then taking the supremum over all
admissible interactions and states, gives
\[
 K\le C+C\sqrt K.
\]
Solving this quadratic inequality for $\sqrt K$ proves a bound
independent of $q$, as required.
\end{proof}

\subsection{Transporting the ground state}

The perturbation is not termwise centered in the sense of
Proposition~\ref{prop:relative-bound}. Spectral transport produces
that property. We give the argument explicitly because it must follow
the perturbed ground line rather than keep the original line fixed.
The smooth frequency filter implements the quasiadiabatic idea
of~\cite{HW} in the exact spectral-flow form of~\cite{BMNS}; the
locality estimates needed here were established
in Section~\ref{sec:locality}.

Let $\varphi$ be real and measurable, $\|\varphi\|_\infty\le1$.
On Fock space set
\[
 V=d\Gamma(T_q(\varphi)),\qquad H_s=H+sV,
 \qquad -1\le s\le1.
\]
The coherent resolution \eqref{eq:toeplitz-resolution} writes $V$
as a bounded-density integral of neutral Hermitian terms
\[
 V_x=\kappa_q\varphi(x)c(k_x)^*c(k_x).
\]
They are rapidly local with uniform bounds by
Lemma~\ref{locality:frame}. The same holds for $H$, using the
rotations of $v_0=e_0\wedge e_1$:
\[
 H=(2q-1)\int_{SU(2)} B(v_0(g))^*B(v_0(g))\,dg.
\]
The measure has uniformly bounded density on the physical sphere.
In particular every $H_s$ has uniform interaction bounds for
$|s|\le1$, even though $\|V\|$ can grow with $N$.

Fix $g_0=1/25$ and work initially on a connected parameter interval
containing zero where the $N$-sector ground state is simple and the
gap above it exceeds $g_0/2$. Such an interval exists at every fixed
flux by finite-dimensional perturbation theory. Choose once and for
all a smooth odd real function $F$ satisfying
\begin{equation}\label{eq:transport-filter}
 F(\omega)=-\frac1\omega\quad (|\omega|\ge g_0/2),
 \qquad F(\omega)=0\quad (|\omega|\le g_0/4).
\end{equation}
It can be chosen with an inverse Fourier kernel having every absolute
polynomial moment. The $1/\omega$ tail is square-integrable and
has integrable derivatives of all positive orders; near time zero
Plancherel gives local integrability, and repeated integration by
parts gives arbitrary decay for large time.

For each term of $V$, define $Y_x(s)$ by applying this frequency
filter with respect to $H_s$. In an eigenbasis of $H_s$,
\begin{equation}\label{eq:transport-matrix-elements}
 \langle a|Y_x(s)|b\rangle
 =F(E_a-E_b)\langle a|V_x|b\rangle.
\end{equation}
Oddness of $F$ makes $Y_x$ anti-Hermitian. Every $Y_x$ is physical,
neutral, and rapidly local with uniform bounds; put
$Y(s)=\int Y_x(s)\,dA(x)$ and solve
\[
 U'(s)=Y(s)U(s),\qquad U(0)=I.
\]
This unitary preserves number. If $P_s$ is the $N$-sector ground
projection, the usual spectral derivative formula and
\eqref{eq:transport-filter} give $P_s'=[Y(s),P_s]$.
Consequently $U(s)P_0U(s)^*=P_s$. Neutrality is important here:
the matrix elements crossing the ground projection stay in the
$N$-particle sector. No assertion about a perturbed Fock-space gap
is needed.

For each $x$ separately, define
\begin{equation}\label{eq:transport-centered-terms}
 \begin{split}
 Q_x(s)&=U(s)^*\bigl(V_x+[H_s,Y_x(s)]\bigr)U(s),\\
 a_x(s)&=\langle\Omega,Q_x(s)\Omega\rangle,
 \qquad A_x(s)=Q_x(s)-a_x(s)I.
 \end{split}
\end{equation}
These are Hermitian neutral physical operators. For an excited
$N$-sector eigenvector $a$ and the ground vector $0$, their
unconjugated off-diagonal matrix element is
\[
 \bigl(1+(E_a-E_0)F(E_a-E_0)\bigr)
       \langle a|V_x|0\rangle=0.
\]
Hermiticity gives the reverse vanishing. Thus each $Q_x(s)$
preserves $\mathbb C\Omega$, and $A_x(s)\Omega=0$ exactly.
The commutator with $H_s$ and the finite-parameter conjugation by
$U(s)$ preserve rapid locality, uniformly on the bootstrap interval.
The scalar satisfies $|a_x(s)|\le\|Q_x(s)\|$; it can be assigned
to the same center. Hence the interaction $A_x(s)$ satisfies
Proposition~\ref{prop:relative-bound} with a uniform size bound.

Let $E_0(s)$ be the $N$-sector ground energy. Differentiating the
pulled-back Hamiltonian and taking its expectation in $\Omega$
gives
\begin{equation}\label{eq:transport-relative-derivative}
 \frac{d}{ds}\bigl(U(s)^*H_sU(s)-E_0(s)I\bigr)
 =\int A_x(s)\,dA(x),
 \qquad
 -CG\le\int A_x(s)\,dA(x)\le CG.
\end{equation}
All differentiation is finite-dimensional; the integral notation
may equally be collected into cells. On the $N$-sector, $G=H$.
Integrating in either parameter direction yields
\begin{equation}\label{eq:final-form-bound}
 U(s)^*H_sU(s)-E_0(s)I\ge(1-C|s|)H
 \qquad\text{on }\bigwedge^N U_q.
\end{equation}
The kernel on this interval is $\mathbb C\Omega$, so the min--max
principle gives a gap of at least $(1-C|s|)g_0$.

Choose $\lambda_*>0$ with $\lambda_*\le1$ and
$C\lambda_*\le1/4$. The improved lower bound $3g_0/4$ prevents
the gap from reaching the bootstrap boundary $g_0/2$ anywhere in
$[-\lambda_*,\lambda_*]$. Indeed a first boundary point would
contradict continuity of the ordered eigenvalues and
\eqref{eq:final-form-bound}. The construction therefore extends
over this entire interval at every sufficiently large flux. Taking
$\Delta_*=g_0/2$ proves Theorem~\ref{thm:main}.
